% Generated draft: requires mathematical and source review.
\documentclass[11pt]{article}
\usepackage[margin=25mm]{geometry}
\usepackage{fontspec}
\setmainfont{Latin Modern Roman}
\usepackage{amsmath,amssymb,mathtools,unicode-math}
\setmathfont{Latin Modern Math}
\usepackage{enumitem,needspace,xurl,hyperref}
\setlength{\parindent}{0pt}
\setlength{\parskip}{5pt}
\setlength{\emergencystretch}{3em}
\newcommand{\N}{\mathbb{N}}
\newcommand{\Z}{\mathbb{Z}}
\newcommand{\Q}{\mathbb{Q}}
\newcommand{\R}{\mathbb{R}}
\newcommand{\C}{\mathbb{C}}
\newcommand{\F}{\mathbb{F}}
\newcommand{\A}{\mathbb{A}}
\newcommand{\PP}{\mathbb{P}}
\newcommand{\E}{\mathbb{E}}
\newcommand{\eps}{\varepsilon}
\newcommand{\dd}{\,\mathrm{d}}
\begin{document}
\section*{UnsolvedMath ID: 30004077}
\textbf{Source problem code:} catalogue entry
\subsection{Definitions and mathematical statement}
Let \((R,\mathfrak m,k)\) be an excellent regular local ring, with residue field \(k=R/\mathfrak m\), and let \((u,y)=(u_1,\ldots,u_d,y_1,\ldots,y_r)\) be a regular system of parameters. Write \(I=(u_1,\ldots,u_d)R\), and let \(\widehat R\) be the \(\mathfrak m\)-adic completion. For a tuple \(a\) in a ring \(S\), write \((a)S\) for its generated ideal. Let \(f\in R\setminus I\).

\textbf{Unresolved definition.} The record calls \(m=m(f)\) a multiplicity defined from the characteristic polyhedron but does not supply that definition or explicitly specify the construction of \(\operatorname{in}_0(f)\). The following precise reconstruction is conditional on interpreting these as the order and initial form of the image of \(f\) modulo \(I\). This interpretation requires source confirmation.

Under this interpretation, put \(\overline R=R/I\), let \(\overline{\mathfrak m}=\mathfrak m/I\), and denote reduction modulo \(I\) by a bar. Define
\[
 m=\max\{n\in\Z_{\geq0}:\overline f\in\overline{\mathfrak m}^{\,n}\}.
\]
The classes of \(\overline y_1,\ldots,\overline y_r\) identify
\[
 \operatorname{gr}_{\overline{\mathfrak m}}(\overline R)
 \cong k[Y_1,\ldots,Y_r].
\]
Let \(F=\operatorname{in}_0(f)\) be the nonzero homogeneous polynomial of degree \(m\) corresponding to the class of \(\overline f\) in \(\overline{\mathfrak m}^{\,m}/\overline{\mathfrak m}^{\,m+1}\). For a \(k\)-linear subspace \(T\) of the degree-one component \(k[Y_1,\ldots,Y_r]_1\), let \(k[T]\) denote the \(k\)-subalgebra generated by \(T\). Assume
\[
 \nexists\,T\subsetneq k[Y_1,\ldots,Y_r]_1
 \quad\text{such that}\quad F\in k[T].
\]
\textbf{Question, conditional on the interpretation above.} For every tuple \(\varphi_1,\ldots,\varphi_r\in I\widehat R\), put \(\widehat y_j=y_j+\varphi_j\). If
\[
 f\in\bigl((\widehat y_1,\ldots,\widehat y_r)\widehat R\bigr)^m,
\]
must there exist \(z_1,\ldots,z_r\in R\) such that \((u_1,\ldots,u_d,z_1,\ldots,z_r)\) is a regular system of parameters and
\[
 (z_1,\ldots,z_r)\widehat R
   =(\widehat y_1,\ldots,\widehat y_r)\widehat R,
 \qquad
 f\in\bigl((z_1,\ldots,z_r)R\bigr)^m?
\]
Here ideal powers are ordinary ideal powers, with zeroth power equal to the ambient ring; no additional assumption that \(f\) is a nonunit has been inserted.

The supplied original transcription places the shifts in \(IR\) and writes the ideal equality over \(R\); the supplied corrected statement instead places the shifts and ideal equality in \(\widehat R\), as used above. The record attributes this change to restoration of completion notation. These differing formulations must remain distinguishable pending human source review.
\subsection{Short English statement}
The question asks whether formal coordinates that make \(f\) belong to the specified power of their ideal can be realized by parameters in \(R\) generating exactly the same ideal after completion. The initial form must require all the indicated linear directions. This draft uses an explicitly conditional interpretation of the source's undefined multiplicity and initial form, and requires human mathematical review.
\subsection{Sources}
\begin{itemize}
\item Local UnsolvedMath registry: supplied complete record for permanent ID \(30004077\), source problem code \texttt{catalogue entry}, in the Oberwolfach Reports Open Problems collection. Project registry paths are \texttt{data/unsolvedmath/problems.json} and the supplemental \texttt{catalog/problems.json}; their contents were not independently inspected for this draft.
\item Supplied source citation: \emph{Singularities and Homological Aspects of Commutative Algebra} (2020), Oberwolfach Reports, DOI \texttt{10.4171/owr/2019/6}, pp.\ 407--464. \url{https://doi.org/10.4171/owr/2019/6}. The record also labels this source ``Oberwolfach Report 6/2019'' and gives proposed year 2019. No individual proposer is supplied.
\item Supplied literature reference, labelled ``Cossart--Schober'': \url{https://arxiv.org/abs/1411.2522}. No title, theorem number, or precise supporting passage is supplied.
\item Literature assessment dated exactly \texttt{2026-08-22}: the supplied curation labels the problem open and states that the exact extension question for excellent rings remains open, while related characteristic-polyhedron results require additional G-ring, Henselian, or similar hypotheses. This is supplied dated curation, not an independently verified literature assessment. The background expressly requests independent verification of time-sensitive status and cited claims.
\item The supplied statement status is \texttt{corrected\_verified}. Its verification note says the official report was checked, the dangling initial-form condition was expanded to the no-proper-linear-subspace hypothesis, and completion notation was restored. This describes supplied verification history; no source visit or independent verification is claimed here.
\end{itemize}
\end{document}
